## Q2 -- does the fixed-k artifact explain ALL Tukey-flagged plots, or just some? (2026-08-20)

Prompted by a direct challenge to the chapter's own logic: the "10 worst plots -> fixed-k
mechanism traced -> therefore all 266 flagged plots are artifact, not ecology" chain was only
directly verified for 5 example plots. This checks the mechanism against every flagged plot
individually, and characterises the ones it does NOT explain.

### Method

Script: `/Users/shreeyagorasia/.claude/jobs/392cb61a/tmp/q2_outlier_mechanism_check.py`. Loads
`models/growth_curve_attribution/data.py::load_filtered_growth_curve_table("4survey")` (raw
per-survey rows, same filtering as production). Refits `y_max_fit` per plot exactly per
Equation eq:plot-ymax-fit (OLS through the origin, `p4`/`p5` held fixed at yield-class values).
Applies the same Tukey 1.5*IQR fence (Equation eq:tukey-fences) used in the chapter.

**Artifact-signature test**: for each flagged plot, compare observed growth (height at last
survey minus height at first survey) against the assigned yield-class benchmark curve's own
predicted growth over the same age span (`top_height95_yldc_predicted`, already in the table).
If observed growth is SLOWER than the benchmark predicts, that is the exact signature of the
traced mechanism (a plot growing slower than its assigned k gets its asymptote inflated to
compensate, since the fit can't lower k to match the slower reality). If observed growth is NOT
slower (matches or exceeds the benchmark), the implausible asymptote is not obviously explained
by this mechanism.

### Result

My own recomputation flags 250/58,112 plots (0.43%) -- close to but not identical to the
chapter's own reported 266 (0.47%); likely a minor methodology difference (exact filtering
population, or how y_max_fit is computed) between this recomputation and the original pipeline.
Directionally consistent, treat the split below as approximate, not pixel-exact.

- **Artifact-consistent** (observed growth slower than assigned benchmark): 179 / 250 (71.6%).
- **NOT artifact-consistent** (observed growth matches or exceeds benchmark): 71 / 250 (28.4%).

The chapter's current claim ("confirming the mechanism traced above is target construction, not
a genuine ecological anomaly") only actually holds for the 71.6% majority. The other 28.4% are
flagged as implausible for a different, untested reason.

### Environmental characterisation of the two subsets

topex, slope, CanopyCover, time_since_thinning, tas_mean, chelsa_bio12_precip_mm pulled for both
subsets, compared directly against each other (Mann-Whitney U):

| Variable | NOT-artifact median | Artifact-consistent median | MWU p (direct comparison) |
|---|---:|---:|---:|
| **topex** | **+2.58** | **-5.48** | **0.0002** |
| time_since_thinning (mean, % ever-thinned) | 2.05 yrs, 46.5% | 0.78 yrs, 35.2% | 0.051 |
| chelsa_bio12_precip_mm | 2180.10 | 2163.50 | 0.94 (no real difference) |
| slope_degrees | 9.93 | 9.02 | not significant (population-level p=0.115) |
| CanopyCover | 0.80 | 0.69 | not significant (population-level p=0.961) |
| tas_mean_4survey | 8.44 | 8.44 | not significant |

### Read

A real, interesting split, not previously in the chapter. NOT-artifact-consistent plots sit in
notably more SHELTERED terrain (positive topex) than artifact-consistent plots (notably more
EXPOSED, negative topex) -- p=0.0002, the strongest signal found. Plausible account: exposed
plots that grow slower than their assigned yield class (plausible given exposure suppresses
growth) get their asymptote artificially inflated by the fixed-k fit -- the already-traced
mechanism. Sheltered plots that still end up flagged are not explained by that route; a sheltered
plot growing genuinely fast (also plausible, given exposure's known negative effect on growth)
would produce a real, biologically-plausible high asymptote that a fixed yield-class benchmark,
calibrated to typical/average exposure, simply underestimates -- worth stating as a plausible
account, not confirmed further here.

Thinning history is a weaker, borderline signal (p=0.051, more and more-recent thinning in the
not-artifact group) -- consistent with thinning genuinely accelerating individual-tree height
growth beyond what an un-thinned yield-class curve assumes, but underpowered at n=71 to treat as
resolved.

Precipitation is NOT a real differentiator between the two flagged subsets (p=0.94) -- both
groups differ similarly from the general population, so precipitation doesn't distinguish WHICH
mechanism explains a given flagged plot, only that flagged plots in general skew slightly drier
than the population (a separate, weaker observation, not pursued further here).

### Addendum -- chronic wind-stress candidate flag (same day)

Prompted by a direction error found in the existing windthrow check
(`TEMP_rq3_outlier_diagnosis_results_2026-08-12.tex`'s "Wind exposure, storm-year signature"
section): it states flagged plots are "slightly MORE sheltered than the population, the opposite
of what storm damage would predict," then reports topex flagged median -4.03 vs. population
+0.74 -- since higher topex = more sheltered (established in the Q1 nonlinearity work earlier
today), -4.03 vs +0.74 is LOWER, i.e. MORE exposed, not more sheltered. The prose contradicts its
own numbers. That aggregate (-4.03) is consistent with this session's own split: it's the
71.6%-majority artifact-consistent group (median -5.48, exposed) dominating the aggregate; the
28.4% not-artifact group is actually sheltered (+2.58).

The current disturbance flags (clearfell: >=25%/70%/70% drops; measurement-inconsistency:
similar) only catch ABRUPT, catastrophic events. They cannot catch CHRONIC wind stress -- a site
growing steadily slower than its yield class assumes with no sudden drop at all -- which would
produce exactly the same "growth deficit vs. benchmark" signature already used to define
"artifact-consistent" above. Built and tested a new, non-abrupt flag directly:
**chronic wind-stress candidate = bottom-quartile topex (population threshold -10.07) AND
growth deficit vs. yield-class benchmark** (same growth-deficit test as above), computed for
every plot in the population, not just the 250 flagged ones
(`/Users/shreeyagorasia/.claude/jobs/392cb61a/tmp/q2_windstress_flag_check.py`).

Because "artifact-consistent" already requires growth deficit by construction, the fair
comparison is against OTHER growth-deficit plots, not the whole population (an unconditional
population comparison overstates the effect, checked and rejected as circular):

| Group | n | % also high-exposure |
|---|---:|---:|
| All growth-deficit plots, population-wide | 20,707 | 29.4% |
| Artifact-consistent (Tukey-flagged) group | 179 | **41.3%** |

Chi-square, artifact-consistent rate vs. the growth-deficit population baseline: chi2=11.54,
p=0.0007. Rate ratio 1.40x -- real, statistically significant, but modest: 58.7% of the
artifact-consistent group is NOT high-exposure, so chronic wind stress does not explain most of
this group either.

### Read (revised)

Chronic wind exposure is a real, partial contributor to WHY some growth-deficit plots become
extreme (Tukey-flagged) outliers specifically, not just ordinary growth-deficit plots -- a
genuine, if modest, enrichment (1.40x, p<0.001) that the current disturbance flags structurally
cannot detect, since they require an abrupt drop. This does not overturn the fixed-k mechanism
account (which remains a real, traceable, purely mechanical effect for at least some of these
plots, per the 5-plot direct trace already in the chapter) -- it adds a second, partial,
non-exclusive contributor specifically within the "artifact-consistent" group, and corrects the
existing windthrow note's own direction error. The 28.4% not-artifact-consistent group's own
sheltered/positive-topex profile is unaffected by this addendum -- that group's story (candidate
for real, unusually fast growth in good conditions) stands as originally found.

### Addendum -- boundary-clustering test (same day), inconclusive, and why

Tried a further test: does clearfell_like or artifact-consistent status respect sub-compartment
management boundaries (whole-unit felling) or cross them irregularly (more consistent with
scattered environmental damage)? Measured, for each flagged plot, what fraction of its own
sub-compartment (scpt) is also flagged.

Result: both groups look similarly scattered -- clearfell_like median 0.6% of same-unit plots
also flagged; artifact-consistent median 0.4%. Neither shows the block-shaped, whole-unit
co-occurrence a complete felling operation would produce.

**This does not mean neither group is management-related.** Felling is not always complete --
thinning (removing some trees, not the whole stand) and partial/salvage fellings are routinely
patchy and selective within a management unit by design, not uniform across it. Low within-unit
clustering is therefore consistent with BOTH a genuine but patchy management intervention and
genuine scattered environmental disturbance -- this test cannot distinguish the two on spatial
pattern alone. Resolving it would need felling-licence records or finer within-unit management
data, neither available here. Reported as a genuinely unresolved limitation, not pushed further.

### Caveats

n=71 for the not-artifact subset -- real but modest statistical power; topex's result is strong
(p=0.0002) and should hold, thinning's (p=0.051) is borderline and should be read cautiously.
Recomputed flagged set (250) is close to but not identical to the chapter's reported 266 -- a
methodology reconciliation, not a full pipeline rerun, would be needed before citing exact
percentages from this note as chapter numbers without re-deriving them from the actual pipeline.
The "sheltered plots grow genuinely fast" explanation is plausible and consistent with the wind/
exposure literature already used elsewhere in this project's background, but not directly tested
here (would need a free-k refit or an independent growth-rate check on just this subset, not run
-- time budget).
